
Aplikácia pozostáva z troch častí: backendu, AI asistenta a mobilnej aplikácie. Backend a asistent bežia v Docker kontajneroch, mobilná aplikácia sa nasadzuje na Android zariadenie.

\subsection*{Predpoklady}

Na hostiteľskom zariadení musí byť nainštalovaný \textbf{Docker} a \textbf{Docker Compose}. Pre buildovanie mobilnej aplikácie je potrebný \textbf{Flutter SDK} verzia 3.24 alebo novšia s nastaveným Android toolchainom. Edge zariadenie aj mobilný telefón musia byť pripojené v rovnakej lokálnej sieti.

\subsection*{Konfigurácia citlivých údajov}

Pred prvým spustením je nutné nastaviť dve hodnoty:

\begin{itemize}
      \item \texttt{JWT\_SECRET} v súbore \texttt{backend/.env} a v \texttt{assistant/docker-compose.yml} \textbf{musí byť rovnaká, dlhá náhodná hodnota}. Tento kľúč podpisuje JWT tokeny a zdieľajú ho obe služby na ich validáciu. Vygenerovať ho je možné napríklad cez \texttt{openssl rand -hex 32}.
      \item \texttt{ANTHROPIC\_API\_KEY} v súbore \texttt{assistant/docker-compose.yml} musí obsahovať vlastný kľúč zo služby \url{https://console.anthropic.com/}, inak AI asistent nebude funkčný.
\end{itemize}

\subsection*{Spustenie backendu a asistenta}

V dvoch oddelených termináloch spustite:

\begin{lstlisting}[language=bash]
cd backend && docker compose up -d --build
cd assistant && docker compose up -d --build
\end{lstlisting}

Po prvom spustení asistenta je potrebné stiahnuť embedovací model:

\begin{lstlisting}[language=bash]
docker exec -it ollama-llm ollama pull nomic-embed-text
\end{lstlisting}

Backend bude dostupný na porte \textbf{8000}, asistent na porte \textbf{8100}. Backend pri prvom štarte automaticky vytvorí databázové tabuľky, založí preddefinované účty a zaregistruje testovaciu linku \emph{Dopravník} s ArUco~ID~1.

\subsection*{Spustenie mobilnej aplikácie}

Pripojte Android zariadenie cez USB s povoleným režimom \emph{USB debugging} a v priečinku \texttt{frontend/} spustite:

\begin{lstlisting}[language=bash]
flutter pub get
flutter run --release
\end{lstlisting}

Po prvom otvorení aplikácie je potrebné v \emph{Nastaveniach} nastaviť IP adresu edge zariadenia pre backend (\texttt{http://<edge-ip>:8000}) aj pre asistenta (\texttt{http://<edge-ip>:8100}) a aplikáciu reštartovať.

\subsection*{Prvé prihlásenie}

Pre prvé prihlásenie sú dostupné preddefinované účty podľa hodnôt v \texttt{backend/.env} (predvolene admin: \texttt{admin@admin.com} / \texttt{admin}). Heslá je odporúčané po prvom prihlásení zmeniť v profile.
